\documentclass[12pt]{article}
\usepackage{graphicx}
\usepackage[margin=2cm, a4paper]{geometry}
\usepackage{setspace}
\usepackage{pdfpages}
\usepackage{float}
\usepackage{amsmath}
\usepackage{fancyvrb}
\usepackage{amssymb}
% \usepackage{minted}
\usepackage{enumitem}
\usepackage{amsthm}
% \usepackage[colorlinks,linkcolor=blue]{hyperref}


% \renewcommand{\contentsname}{Contents}
% \renewcommand{\figurename}{Figure}
% \renewcommand{\tablename}{Table}
% \hypersetup{
    % colorlinks=true,
    % linkcolor=black,
    % filecolor=magenta,      
    % urlcolor=blue,
% }
\newcommand{\mytitle}{Abstract Algebra II - Homework 3}
\newcommand{\myauthor}{B13902022 Yu-Chi,Lai}

\usepackage{fancyhdr}
\pagestyle{fancy}
\fancyhead{}
\fancyhead[L]{\mytitle}
\fancyhead[R]{\myauthor}

\usepackage[english]{babel}
\newtheorem{theorem}{Theorem}[section]
\newtheorem{lemma}[theorem]{Lemma}

\theoremstyle{definition}
\newtheorem{definition}{Definition}[section]
\theoremstyle{remark}
\newtheorem*{remark}{Remark}
\onehalfspacing

\title{\mytitle}
\author{\textbf{\myauthor}}
\date{}

\begin{document}

\maketitle

\setcounter{section}{1}
\section*{1}
\subsection*{(a)}
Let $f(x)=c_mx^m+c_{m-1}x^{m-1}+\dots+c_1x+c_0, c_i\in K\forall i\in\{0,1,\dots,m\}$. Let $\sigma\in G$ and $a_i\in X$, we have: (Clearly all elements in $G$ fix $K$ and $f(a_i)=0$ and $\sigma(0)=0$)
\begin{align*}
\sigma(f(a_i))&=\sigma(c_ma_i^m)+\sigma(c_{m-1}a_i^{m-1})+\dots+\sigma(c_1a_i)+\sigma(c_0) \\
&=c_m(\sigma(a_i))^m+c_{m-1}(\sigma(a_i))^{m-1}+\dots+c_0 \\
&=f(\sigma(a_i))=0
\end{align*}
Thus, $\sigma(a_i)$ is also a root of $f$, so $\sigma(a_i)\in X$, since $\sigma$ is a bijection on the field $L$, its restriction to the finite set $X$ is a bijection from $X$ to itself. Hence, every $\sigma \in G$ defines a permutation of $X$.

Define $\phi:G\rightarrow \mathrm{Perm(X)},\ \sigma\in G$ by $\phi(\sigma)=\sigma|_{X}$ (the restriction of $\sigma$ to the set $X$). For any $\sigma, \tau\in G$ and any $a\in X$, we have:
\[\phi(\sigma\circ\tau)(a)=(\sigma\circ\tau)(a)=\sigma(\tau(a))=\phi(\sigma)(\phi(\tau)(a))\]
It's obviously a homomorphism. We then try to prove the map is injective, i.e., the kernel of $\phi$ is the identity automorphism. By definition, the splitting field $L$ is the smallest field containing $K$ and all the roots of $f$, therefore, $L$ is generated by the roots over $K$, in other words, $L=K(a_1,a_2,\dots,a_n)$.

Suppose $g\in \ker{(\phi)}$, any field automorphism in $G$ is completely determined by its action on the generators of the field extension. Since $g$ fixes $K$ and fixes every generator $a_i$, it must fix every element of $L$, thus, $g$ is the identity map. Because $\ker{\phi}={id_L}$, the map $\phi$ is injective.
\subsection*{(b)}

\begin{theorem}[Isomorphism extension theorem]
Given any field $F$, an algebraic extension field $E$ of $F$ and an isomorphism $\varphi$ mapping $F$ onto a field $F'$ then $\varphi$ can be extended to an isomorphism $\tau$ mapping $E$ onto an algebraic extension $E'$ of $F'$ (a subfield of the algebraic closure of $F'$). 
\end{theorem}
\begin{theorem}
Let $F$ be a field of characteristic $0$ and let $f(x)\in F[x]$, and let $G$ be the Galois group of $f(x)$ over $F$. If $f(x)$ is irreducible, $G$ acts transitively on the roots of $f(x)$.
\end{theorem}
\begin{proof}
Because $f$ is an irreducible polynomial over $K$, for any two roots $a_i,a_j\in X$, there exists a field isomorphism between the simple extensions:
$$\theta:K(a_i)\rightarrow K(a_j)$$

that fixes $K$ and maps $a_i$ to $a_j$.
Because $L$ is the splitting field of $f$ over both $K(a_i)$ and $K(a_j)$, this isomorphism $\theta$ can be extended to an automorphism $\sigma$ of the entire splitting field $L$ (by the Isomorphism Extension Theorem).

Therefore, there exists some $\sigma \in G$ such that $\sigma(a_i)=a_j$. This means the action of $G$ on the set of roots $X$ is transitive (there is only one orbit).
\end{proof}

By the orbit-stabilizer theorem, $G$ acts on set $X$, the size of the orbit of an element multiplied by the size of its stabilizer subgroup equals the size of the group:
\[|G|=|\mathrm{Orbit}(a_i)||\mathrm{Stab}_G(a_i)|\]
Since the action is transitive, $|\mathrm{Orbit}(a_i)|=|X|=n$, hence $n$ divides $|G|$.







\setcounter{section}{2}
\section*{2}
\begin{theorem}
Symmetric group $S_p$ (where $p$ is a prime) can be generated by a transposition and a p-cycle.
\end{theorem}
\begin{proof}
Consider any transposition $\tau=(ab)$, and the p-cycle $\sigma=(12\dots p)$. Because $\langle \sigma\rangle$ is the subgroup of $S_p$ consists of p-cycles, and every element has order $p$ ($p$ is prime) and can be the generator. Hence, there exists integer $k$ such that $\sigma^k$ sends $a$ to $b$ and is a p-cycle.
\end{proof}

\end{document}